\chapter{Business Model and Pricing}\label{ch:business}

\section{Who pays for what}\label{sec:pay}

The retailer pays, and it pays per order served. The shopper uses the widget for free. This follows from where the benefit lands: the saving that ClosetTwin promises is the cost of the returns that no longer happen, and that saving grows with the number of orders the retailer processes. A fee per order therefore moves with the value delivered, it lets a retailer start small, and it avoids a fixed license that a retailer would have to justify before it has seen any effect. The only other revenue in the model is the fixed fee for the pilot. The data feed described in Chapter~\ref{ch:product} earns nothing in this plan.

\section{Value-based price}\label{sec:pricing-value}

The price is set against the value, and the value is the return cost that the tool avoids. We follow the logic of Gustafsson et al., who compute what a fitting technology may cost per unit before it stops paying for itself. In their case study of a footwear retailer, with a prime cost assumed at EUR~30, the allowed cost is EUR~1.1 per pair if the technology verifies fit in 40\% of cases, EUR~1.6 at 60\% and EUR~2.1 at 80\% \cite{gustafsson2021}. Those figures belong to one retailer and one category, so we rebuild the calculation with our own inputs; formula~(\ref{eq:breakeven}) in Chapter~\ref{ch:finance} gives it symbol by symbol. The cost of a return is \sourced{17\%} of the prime cost \cite{gustafsson2021}. With an \assumed{\$80} order and a prime cost that we assume to be \assumed{50\%} of it, the prime cost is \derived{\$40} and a return costs \derived{\$6.80}. The return rate is \sourced{24.4\%} \cite{coresight2023}, and we assume that \assumed{40\%} of returns are fit-related. Together this is \derived{\$0.66} of fit-related return cost per order. If ClosetTwin cuts these returns by \assumed{15\%}, the saving is \derived{\$0.0996}, about \$0.10 per order. This is the most the tool could cost and still break even for the retailer.

The reduction of 15\% is deliberately modest. Gustafsson et al. test reductions between 40\% and 80\% \cite{gustafsson2021}, but the largest field study that we have found shows size-finder users slightly \emph{more} likely to return an item \cite{patel2025}, so we start well below the scenarios and let the pilot decide. The prices are set as a fraction of this value. The Growth tier charges \assumed{\$0.04} per order, which is 40\% of the break-even value, and the Enterprise tier \assumed{\$0.03}, which is 30\%. At these prices a retailer breaks even if the true reduction of fit-related returns is at least \derived{6.03\%} (Growth) or \derived{4.52\%} (Enterprise), and the retailer keeps the rest of the saving. Under our mix of 60\% Growth and 40\% Enterprise clients, the blended price is \derived{\$0.036} per order, about \ils{0.13} at 3.70 shekels to the dollar.

\section{The tiers}\label{sec:tiers}

Table~\ref{tab:tiers} shows the three tiers. The pilot is a fixed fee that covers integration and a four-week test; the amount depends on the market, because a mid-size Israeli client, as the next section shows, pays about \$3,600 a year and cannot be asked for \$5,000 before it has seen any result. The Growth and Enterprise tiers charge per order served, and Enterprise carries a volume discount. Low-touch pilots through an app store, which the plan uses only in Europe, are free. All prices are in US dollars.

\begin{table}
\caption{Pricing tiers (all values are assumptions)}\label{tab:tiers}
\begin{tabularx}{\textwidth}{lXRR}
\toprule
\ctheaderrow\cthead{Tier} & \cthead{What it covers} & \cthead{Israel} & \cthead{Europe}\\
\midrule
Pilot & integration and a four-week randomized test & \assumed{\$2,500} fixed & \assumed{\$5,000} fixed\\
Low-touch pilot & app-store trial, no direct integration & not offered & \assumed{\$0}\\
Growth & per order served; 60\% of paying clients & \assumed{\$0.04} & \assumed{\$0.04}\\
Enterprise & per order served; 40\% of paying clients & \assumed{\$0.03} & \assumed{\$0.03}\\
\bottomrule
\end{tabularx}
\end{table}

\section{Reference points, not benchmarks}\label{sec:price-ref}

The direct enterprise vendors publish no prices (Table~\ref{tab:vendors}), so our prices cannot be benchmarked against them. The only public prices we found are for a scan API and for an app for small merchants, and neither is comparable. 3DLOOK lists its FitXpress API from \$1,000 a month with 500 requests included and \$2.00 for each additional request, and its Mobile Tailor plans from \$499 a month for up to 100 scans and \$999 for up to 500 scans \cite{threedlook_pricing}. The first is a health and fitness scan API and the second a made-to-measure product, so neither is a size recommendation for retail. Popsize's listing on the Shopify App Store states a free plan and a free trial and shows no paid tier \cite{popsize_shopify}. We show these as reference points only. They say what a scan costs when bought from a vendor and what a small merchant is offered for free, and neither says what a large retailer would pay for a result.

\section{Unit economics and the scan-cost condition}\label{sec:unit-economics}

Gross margin is high only if a scan is cheap. We assume that a scan costs us \assumed{\$0.03} and that \assumed{25\%} of orders include one, which gives a cost to serve of \derived{\$0.0075} per order and a gross margin of \derived{79\%} at the blended price. At 3DLOOK's list price per scan, \$999 for 500 scans or about \$2.00 each \cite{threedlook_pricing}, a scan would cost about \derived{67 times} our assumption, and every scan would cost far more than the order pays. The business model therefore depends on in-house scan capacity, and Chapter~\ref{ch:risks} treats a scan cost above the model as a risk of its own.

Israel adds a second honest limit. The largest Israeli fashion retailer in our bands has 600,000 orders a year, so the average Israeli paying client, with 30\% large and 70\% mid-size retailers, sends about \derived{250,000} orders a year and pays about \derived{\$9,000} a year at the blended price. The price per order is right for the client, but the ticket is small, and the whole Israeli SAM is worth about \$198,000 a year (Chapter~\ref{ch:market}). A European client with the assumed 1,000,000 orders a year pays \derived{\$36,000}, which is why the plan depends on Europe. For comparison, the saving of such a European client at the base case is about \$99,600 a year against a Growth fee of \$40,000, and it is real only if the pilot confirms the effect.

\ctchapterend
